\documentclass[preprint,12pt]{elsarticle}

\usepackage[utf8]{inputenc}
\usepackage{amsmath,amssymb,amsfonts,amsthm}
\usepackage{graphicx}
\usepackage{booktabs}
\usepackage{multirow}
\usepackage{array}
\usepackage{algorithm}
\usepackage{algpseudocode}
\usepackage{hyperref}
\usepackage{xcolor}
\usepackage{subcaption}
\usepackage{microtype}
\usepackage{geometry}

\geometry{
  a4paper,
  top=25mm,
  bottom=25mm,
  left=25mm,
  right=25mm
}

\journal{Computers \& Security}

\begin{document}

\begin{frontmatter}

\title{A Hybrid Deep Learning, Quantum Variational, and Density-Based Latent Clustering Architecture for Zero-Day Network Intrusion Detection and Forensic Explainability}

\author[inst1]{Author Name\corref{cor1}}
\ead{author@institution.edu}
\author[inst1]{Co-Author Name}
\ead{coauthor@institution.edu}

\affiliation[inst1]{organization={Department of Computer Science and Cybersecurity},
            addressline={University / Research Center}, 
            city={City},
            postcode={PostalCode},
            state={State},
            country={Country}}

\cortext[cor1]{Corresponding author}

\begin{abstract}
Zero-day network intrusions represent an asymmetric operational challenge for modern Security Operations Centers (SOCs). Classical signature-driven Intrusion Detection Systems (IDS) fail against previously unobserved malicious payloads, while monolithic Deep Learning (DL) architectures frequently suffer from high false-alarm rates, opaque decision surfaces, and severe compression of novel structural anomalies in Euclidean latent space. To address these fundamental limitations, this paper proposes an end-to-end, multi-stage cascaded architecture integrating a One-Dimensional Convolutional Neural Network (1D-CNN), a 6-qubit Quantum Siamese Neural Network (QSNN), and a latent Density-Based Spatial Clustering of Applications with Noise (DBSCAN) triage engine, coupled with an Explainable Artificial Intelligence (XAI) attribution module based on Shapley Additive exPlanations (SHAP). In the proposed pipeline, high-throughput classical 1D-CNN screening clears 47.12\% of standard baseline traffic instantly, forwarding suspicious flows to a parameterized quantum circuit that exploits Hilbert space embedding via data re-uploading and Triplet Margin Metric Learning with an intra-class spread penalty. Borderline ambiguous flows landing in a quantum uncertainty band ($P_{\text{qsnn}} \in [0.45, 0.60)$) are escalated to DBSCAN operating across 10 orthogonal Principal Component Analysis (PCA) dimensions of the CNN latent manifold, anchored by 22 normal and 14 attack centroids. Rigorous evaluation on the complete 82,332-flow test cohort of the UNSW-NB15 benchmark—with the Worms attack family strictly held out during training—demonstrates an overall system accuracy of 90.31\%, a detection precision of 92.90\%, and an attack recall of 89.21\%. In the uncertainty zone, DBSCAN successfully isolates 15 zero-day outliers with 100.0\% attack purity while salvaging 9 benign novel connections from false alarm generation. Furthermore, threshold-free novelty analysis confirms that quantum representations preserve structural novelty significantly better than classical embeddings, with a Quantum Autoencoder (QAE) reconstruction error achieving an AUROC of 0.7735 compared to $\sim$0.55 for classical latent distances. Finally, an interactive enterprise SOC dashboard provides game-theoretic feature attribution with non-negative visual domains, delivering transparent, verifiable forensic telemetry for mission-critical cyber defense operations.
\end{abstract}

\begin{highlights}
\item Multi-stage cascaded intrusion detection architecture fusing 1D-CNN, 6-qubit QSNN, and DBSCAN.
\item Metric learning in Hilbert space via Triplet Margin Loss and data re-uploading achieving an inter-to-intra cluster separation ratio of 3.7551.
\item Geometric isolation of zero-day outliers with 100.0\% attack purity and automated mitigation of false alarms.
\item Empirical evidence that quantum representations preserve structural zero-day novelty (AUROC 0.7735) where classical latent distances compress it away.
\item Game-theoretic forensic explainability via SHAP with non-negative visual attribution for Security Operations Centers.
\end{highlights}

\begin{keywords}
Zero-Day Intrusion Detection \sep Quantum Machine Learning \sep Quantum Siamese Neural Networks \sep DBSCAN Clustering \sep Explainable AI \sep Network Forensics
\end{keywords}

\end{frontmatter}

\section{Introduction}
\label{sec:intro}

The proliferation of sophisticated Advanced Persistent Threats (APTs), polymorphic malware, and automated exploitation frameworks has fundamentally altered the threat landscape of modern enterprise networks \cite{sommer2010outside, buzakova2020survey}. Traditional Network Intrusion Detection Systems (NIDS) primarily deploy signature-matching engines (e.g., Snort, Suricata, Zeek) that compare observed packet headers and payload byte sequences against curated repositories of known malicious indicators \cite{roesch1999snort}. While computationally efficient and capable of processing line-rate multi-gigabit traffic, signature-based approaches possess an inherent, fatal architectural vulnerability: they are categorically blind to zero-day vulnerabilities, novel protocol mutations, and evasion techniques that diverge from historical catalogs \cite{ahmad2021network}.

To circumvent the rigid limitations of signature engines, the cybersecurity research community has invested heavily in Machine Learning (ML) and Deep Learning (DL) paradigms \cite{shone2018deep, ferrag2020deep}. Supervised classifiers—such as Random Forests, Support Vector Machines (SVMs), Deep Neural Networks (DNNs), and Convolutional Neural Networks (CNNs)—treat intrusion detection as a statistical pattern recognition task \cite{vinayakumar2019deep}. By training on high-dimensional flow records comprising transmission durations, inter-arrival intervals, byte volumes, and flag states, these models can generalize beyond static signatures to identify broad behavioral patterns of malicious activity. 

However, deploying monolithic deep learning classifiers into mission-critical Security Operations Centers (SOCs) reveals four critical architectural bottlenecks:
\begin{enumerate}
    \item \textbf{The False Alarm Tsunami and Alert Fatigue:} Supervised deep networks optimize global empirical risk over closed-world training distributions. When subjected to enterprise traffic containing legitimate but unobserved protocol updates or bursty application behaviors, classical models exhibit severe overconfidence, emitting thousands of false-positive alarms daily that overwhelm human analysts \cite{axelsson2000base}.
    \item \textbf{Latent Space Compression of Novel Outliers:} In classical feed-forward neural networks, successive pooling and nonlinear activation operations compress input vectors into Euclidean manifolds optimized exclusively to separate observed classes. Consequently, genuine zero-day attacks often project dangerously close to familiar attack centroids or benign clusters, obliterating the fine-grained topological signatures necessary to flag them as unobserved phenomena \cite{chandola2009anomaly}.
    \item \textbf{Computational Throughput vs. Complex Inference:} Complex anomaly detectors and non-parametric clustering algorithms exhibit quadratic time complexity $\mathcal{O}(N^2)$, rendering them intractable for continuous high-throughput line-rate screening \cite{ester1996density}. Conversely, lightweight linear models lack the representational capacity to resolve sophisticated polymorphic intrusions.
    \item \textbf{The Opacity of "Black-Box" Decisions:} High-dimensional deep neural networks do not provide human-interpretable justifications for their inferences. When an alert fires, SOC analysts are left without actionable insight into which specific packet attributes (e.g., Time-to-Live discrepancies, window advertisements, or byte asymmetry) triggered the automated firewall drop rule or quarantine directive \cite{guidotti2018survey}.
\end{enumerate}

To simultaneously overcome these challenges, this paper presents a novel, hierarchical cyber defense architecture that synergistically combines classical deep learning, quantum variational metric learning, unsupervised density-based clustering, and game-theoretic explainability into a unified, four-stage cascaded pipeline. 

\subsection{Key Contributions}
The primary contributions of this work are summarized as follows:
\begin{itemize}
    \item \textbf{Hierarchical Cascaded Pipeline Architecture:} We design and implement an end-to-end multi-stage intrusion triage architecture. A classical 1D-CNN serves as a high-speed line-rate screener that instantly clears 47.12\% of standard baseline traffic, shielding computationally intensive downstream stages from unnecessary overhead.
    \item \textbf{Metric Learning in Quantum Hilbert Space:} We formulate and implement a 6-qubit Quantum Siamese Neural Network (QSNN) utilizing Parameterized Quantum Circuits (PQCs) with data re-uploading and circular CNOT entanglement. Trained via a custom Triplet Margin Loss regularized by an intra-class spread penalty, the QSNN maps non-linear flow correlations into an expanded Hilbert state space, achieving an inter-to-intra cluster separation ratio of 3.7551 and an attack recall of 99.76\%.
    \item \textbf{Geometric Latent Triage via DBSCAN:} We introduce an automated escalation mechanism where borderline flows exhibiting quantum ambiguity ($P_{\text{qsnn}} \in [0.45, 0.60)$) are routed to a 10-dimensional PCA latent space populated by 22 normal and 14 attack centroids extracted via DBSCAN. This provides geometric isolation of true zero-day outliers with 100.0\% attack purity while recovering benign novel flows to prevent alert fatigue.
    \item \textbf{Empirical Proof of Quantum Novelty Preservation:} Through rigorous leave-one-attack-out validation on the UNSW-NB15 benchmark with the Worms family strictly held out during training, we provide empirical evidence that quantum representations preserve structural zero-day novelty (QAE reconstruction AUROC of 0.7735) where classical latent distance metrics degrade to near-random performance ($\sim$0.55 AUROC).
    \item \textbf{Game-Theoretic Forensic Explainability and Operational SOC Telemetry:} We integrate a customized KernelSHAP module that decomposes neural classifications into strictly non-negative feature attribution impact magnitudes ($|\text{SHAP}| \ge 0$) paired with directional risk indicators and automated plain-text incident narratives, deployed within a live enterprise SOC dashboard.
\end{itemize}

\subsection{Paper Organization}
The remainder of this paper is structured as follows. Section \ref{sec:related} reviews related work in classical NIDS, Quantum Machine Learning (QML), and Explainable AI. Section \ref{sec:formulation} establishes the formal threat model, mathematical problem formulation, and dataset preparation protocol. Section \ref{sec:methodology} provides a comprehensive exposition of the proposed multi-stage hybrid methodology across all constituent modules. Section \ref{sec:experiments} details the experimental setup, simulation environment, and hyperparameter specifications. Section \ref{sec:results} presents extensive empirical results, ablation studies, quantum comparative analyses, and stress-testing evaluations. Section \ref{sec:xai} examines forensic explainability case studies. Section \ref{sec:discussion} discusses operational feasibility, computational complexity, and security implications. Finally, Section \ref{sec:conclusion} concludes the paper and identifies future research trajectories.

---

\section{Related Work and Theoretical Foundations}
\label{sec:related}

\subsection{Deep Learning for Network Intrusion Detection}
Over the past decade, deep learning architectures have superseded shallow machine learning algorithms in network anomaly detection due to their hierarchical representation learning capabilities \cite{al2018intrusion, khraisat2019survey}. Vinayakumar et al. \cite{vinayakumar2019deep} demonstrated that Deep Neural Networks (DNNs) consistently outperform shallow models such as Naive Bayes, Decision Trees, and Random Forests across the KDD Cup 99, NSL-KDD, and UNSW-NB15 benchmarks. Shone et al. \cite{shone2018deep} proposed non-symmetric deep autoencoders for unsupervised feature extraction, demonstrating reduced training times and competitive classification accuracy.

Convolutional Neural Networks (CNNs), initially popularized in computer vision, have been successfully adapted for tabular and sequential network flow data \cite{wang2017malware}. By treating consecutive flow attributes as one-dimensional spatial signals, 1D-CNNs leverage local receptive fields and weight sharing to capture inter-feature correlations (such as cross-packet byte-to-duration ratios and TCP flag transitions) while maintaining invariance to minor feature perturbations. However, when deployed as monolithic classifiers, supervised CNNs suffer from closed-world distributional bias: they assign test samples to the nearest known class probability simplex, rendering them incapable of isolating previously unobserved attack vectors without significant retraining \cite{bendale2016towards}.

\subsection{Quantum Machine Learning in Cybersecurity}
Quantum Machine Learning (QML) investigates the intersection of quantum information theory and statistical learning \cite{biamonte2017quantum, cerezo2021variational}. In the current Noisy Intermediate-Scale Quantum (NISQ) era, Parameterized Quantum Circuits (PQCs)—also designated as Variational Quantum Algorithms (VQAs)—represent the leading paradigm for realizing quantum utility on NISQ-era processors and high-performance state-vector simulators \cite{preskill2018quantum, mcclean2016theory}.

In a PQC, classical feature vectors $\mathbf{x} \in \mathbb{R}^D$ are mapped into the state space of an $n$-qubit register via unitary state preparation operators $U_\phi(\mathbf{x})$, creating an $n$-qubit quantum state:
\begin{equation}
    |\psi(\mathbf{x})\rangle = U_\phi(\mathbf{x}) |0\rangle^{\otimes n}
\end{equation}
The dimensionality of this state space scales exponentially as $\mathcal{H} \cong \mathbb{C}^{2^n}$. A parameterized ansatz $V(\boldsymbol{\theta})$ comprising single-qubit rotations and multi-qubit entangling operations (e.g., CNOT gates) transforms the state:
\begin{equation}
    |\psi(\mathbf{x}, \boldsymbol{\theta})\rangle = V(\boldsymbol{\theta}) U_\phi(\mathbf{x}) |0\rangle^{\otimes n}
\end{equation}
Expectation values of Hermitian observables $\hat{\mathcal{O}}$ are measured to yield real-valued predictions:
\begin{equation}
    \langle \hat{\mathcal{O}} \rangle = \langle \psi(\mathbf{x}, \boldsymbol{\theta}) | \hat{\mathcal{O}} | \psi(\mathbf{x}, \boldsymbol{\theta}) \rangle
\end{equation}
Schuld et al. \cite{schuld2019quantum, schuld2021effect} established that PQCs can be interpreted as quantum kernel estimators and Fourier series approximators in Hilbert space. Havl{\'\i}cek et al. \cite{havlicek2019supervised} demonstrated that quantum circuits can generate feature maps that are classically intractable to simulate, offering potential computational advantages for complex classification tasks.

In cybersecurity, early exploratory applications of QML have focused on Quantum Support Vector Machines (QSVMs) \cite{rebentrost2014quantum}. However, standard QSVM implementations require the evaluation of an $\mathcal{O}(N^2)$ quantum kernel matrix. For realistic enterprise cohorts (e.g., $N = 175,341$ samples in the UNSW-NB15 training set), computing the full Gram matrix on quantum simulators would require approximately $1.53 \times 10^{10}$ circuit executions—a computational barrier requiring years of runtime. In contrast, variational Quantum Neural Networks (QNNs) optimize a fixed parameter vector $\boldsymbol{\theta}$, enabling linear $\mathcal{O}(N)$ training complexity and $\mathcal{O}(1)$ constant-time query inference.

\subsection{Unsupervised Clustering and Zero-Day Isolation}
To address novel intrusions, unsupervised anomaly detection algorithms identify patterns that deviate statistically from baseline normal profiles \cite{ahmed2016survey}. Density-Based Spatial Clustering of Applications with Noise (DBSCAN), introduced by Ester et al. \cite{ester1996density}, is uniquely suited for cybersecurity applications because it discovers clusters of arbitrary geometric topologies and explicitly isolates low-density outliers as noise ($\text{label} = -1$). Unlike $k$-means, DBSCAN does not assume spherical cluster geometry and does not require pre-specifying the number of clusters $k$.

However, executing DBSCAN directly on raw, high-dimensional packet headers presents two critical challenges: the distance concentration phenomenon (where Euclidean distances in high dimensions become indistinguishable) and heavy computational overhead on large datasets \cite{aggarwal2001surprising}. By projecting deep latent embeddings from a trained neural network into an orthogonal PCA subspace, our proposed architecture leverages DBSCAN specifically on ambiguous flows, eliminating both distance collapse and computational saturation.

\subsection{Explainable Artificial Intelligence (XAI) in Threat Triage}
As machine learning models are integrated into security automation pipelines, explainability has emerged as a fundamental operational requirement \cite{arrieta2020explainable, tjoa2020survey}. Without transparent attribution, security analysts cannot distinguish between genuine malicious indicators and spurious statistical artifacts (e.g., ephemeral port assignments or IP address biases).

Lundberg and Lee \cite{lundberg2017unified} proposed SHAP (SHapley Additive exPlanations), unifying cooperative game theory with local surrogate modeling. Rooted in axiomatic Shapley value formulations \cite{shapley1953value}, SHAP computes the marginal contribution of each feature across all possible subset coalitions. In network forensics, SHAP attributes positive and negative forces that drive an anomaly score, allowing analysts to perform instant root-cause verification \cite{wang2020interpreting}.

---

\section{Threat Model and Problem Formulation}
\label{sec:formulation}

\subsection{Threat Model and Adversarial Assumptions}
We assume an operational enterprise environment monitored by perimeter and internal network sensors capturing bidirectional TCP/IP flow metadata. The threat model incorporates the following adversarial capabilities:
\begin{enumerate}
    \item \textbf{Polymorphic and Evasive Payloads:} Adversaries deploy encryption, payload padding, and timing jitter to defeat byte-level signature detection systems.
    \item \textbf{Exploitation of Zero-Day Vulnerabilities:} Attackers possess functional zero-day exploits targeting previously unpatched software components. These attacks exhibit no historical signature and diverge from known exploit families.
    \item \textbf{Distributional Boundary Mimicry:} Sophisticated threat actors intentionally craft malicious connections to mimic normal operational baseline attributes (e.g., standard packet lengths, typical port usage, and realistic connection durations) in order to land near the decision boundaries of automated screeners.
\end{enumerate}

\subsection{Mathematical Problem Formulation}
Let $\mathcal{X} \subset \mathbb{R}^D$ represent the continuous and categorical feature space of network connection flows, and let $\mathcal{Y} = \{0, 1\}$ denote the binary ground-truth label space, where $y = 0$ corresponds to legitimate normal traffic and $y = 1$ corresponds to a malicious intrusion. 

The global intrusion set $\mathcal{Y}_{\text{attack}}$ is partitioned into two disjoint subsets:
\begin{equation}
    \mathcal{Y}_{\text{attack}} = \mathcal{A}_{\text{known}} \cup \mathcal{A}_{\text{zero-day}}, \quad \mathcal{A}_{\text{known}} \cap \mathcal{A}_{\text{zero-day}} = \emptyset
\end{equation}
During the supervised training phase, the learning system has access only to labeled instances from the training distribution:
\begin{equation}
    \mathcal{D}_{\text{train}} = \{(\mathbf{x}_i, y_i) \mid \mathbf{x}_i \in \mathcal{X}, y_i \in \{0\} \cup \mathcal{A}_{\text{known}}\}_{i=1}^{N_{\text{train}}}
\end{equation}
The zero-day attack family $\mathcal{A}_{\text{zero-day}}$ is strictly withheld from the training phase:
\begin{equation}
    \forall (\mathbf{x}, y) \in \mathcal{D}_{\text{train}}, \quad y \notin \mathcal{A}_{\text{zero-day}}
\end{equation}
The evaluation phase exposes the system to the test distribution $\mathcal{D}_{\text{test}}$, containing both known traffic and unobserved zero-day flows:
\begin{equation}
    \mathcal{D}_{\text{test}} = \mathcal{D}_{\text{test}}^{\text{known}} \cup \{(\mathbf{x}_j, y_j) \mid y_j \in \mathcal{A}_{\text{zero-day}}\}_{j=1}^{N_{\text{zero-day}}}
\end{equation}

The ultimate objective of the cascaded intrusion detection system is to learn a multi-stage mapping $\mathcal{F}: \mathcal{X} \to \mathcal{C}$ that assigns each incoming connection $\mathbf{x}$ to one of four discrete operational operational categories:
\begin{equation}
    \mathcal{C} = \{\text{Normal Traffic}, \text{Malicious Intrusion}, \text{Benign Novel Behavior}, \text{Zero-Day Outlier}\}
\end{equation}
while satisfying the following performance constraints:
\begin{align}
    \text{Throughput}(\mathcal{F}) &\ge \text{LineRateThreshold} \\
    \text{FPR}(\mathcal{F}) &\le \alpha_{\text{operational}} \\
    \text{Precision}(\mathcal{C} = \text{Zero-Day Outlier}) &\to 1.0
\end{align}

\subsection{Benchmark Dataset Protocol (UNSW-NB15)}
To evaluate the proposed architecture, we utilize the widely recognized UNSW-NB15 benchmark dataset \cite{moustafa2015unsw}, constructed in the Cyber Range Lab of the Australian Centre for Cyber Security (ACCS) using an IXIA PerfectStorm tool to simulate realistic modern enterprise network traffic. 

The comprehensive dataset comprises 257,543 flow records split into an official training set of 175,341 records and an independent test set of 82,332 records across 49 features. The dataset encompasses nine distinct attack categories: \textit{Fuzzers, Analysis, Backdoors, DoS, Exploits, Generic, Reconnaissance, Shellcode,} and \textit{Worms}.

\begin{table}[t]
\centering
\caption{Class and attack category distribution across the UNSW-NB15 benchmark with the Worms family strictly isolated as a zero-day holdout.}
\label{tab:dataset_distribution}
\small
\begin{tabular}{lrrcc}
\toprule
\textbf{Traffic Category} & \textbf{Training Flows} & \textbf{Testing Flows} & \textbf{Role in Pipeline} & \textbf{Attack Family ID} \\
\midrule
Normal (Benign)        & 56,000  & 37,000 & Baseline Profile & 0 \\
Generic                & 40,000  & 18,871 & Known Attack     & 1 \\
Exploits               & 33,393  & 11,132 & Known Attack     & 2 \\
Fuzzers                & 18,184  & 6,062  & Known Attack     & 3 \\
DoS                    & 12,264  & 4,089  & Known Attack     & 4 \\
Reconnaissance         & 10,491  & 3,496  & Known Attack     & 5 \\
Analysis               & 2,000   & 677    & Known Attack     & 6 \\
Backdoor               & 1,746   & 583    & Known Attack     & 7 \\
Shellcode              & 1,133   & 378    & Known Attack     & 8 \\
\midrule
\textbf{Worms (Zero-Day)} & \textbf{130 (Held Out)} & \textbf{44} & \textbf{Zero-Day Target} & \textbf{9} \\
\midrule
\textbf{Total Cohort}  & \textbf{175,211} & \textbf{82,332} & -- & -- \\
\bottomrule
\end{tabular}
\end{table}

\textbf{Zero-Day Holdout Protocol:} In our experimental design, all 130 instances of the \textbf{Worms} attack family present in the raw training dataset were systematically extracted and sequestered into an isolated holdout artifact (\texttt{zero\_day\_worms.csv}), leaving $N_{\text{train}} = 175,211$ training flows. The models were trained strictly without any exposure to Worms flows, establishing an uncompromised zero-day testing protocol. The distribution of classes is detailed in Table \ref{tab:dataset_distribution}.

---

\section{System Architecture and Proposed Methodology}
\label{sec:methodology}

The proposed architecture organizes network intrusion detection into a hierarchical four-stage pipeline supported by two cross-cutting modules for feature engineering and explainability. A high-level system schematic is presented in Figure \ref{fig:architecture}.

\begin{figure*}[t]
\centering
\includegraphics[width=\textwidth]{architecture_diagram.png}
\caption{End-to-End System Architecture: Hierarchical cascade coupling (1) Data Preprocessing, (2) 1D-CNN Latent Feature Extraction, (3) 6-Qubit QSNN Verification in Hilbert Space, (4) DBSCAN Latent Geometric Triage, (5) SHAP Forensic Explainability, and (6) Security Operations Center (SOC) Dashboard Orchestration.}
\label{fig:architecture}
\end{figure*}

\subsection{Module 1: Data Preprocessing and Feature Transformation}
Raw network packet attributes exhibit severe scale imbalances, heavy-tailed skewness, and mixed data types. The preprocessing pipeline executes four deterministic transformations:

\subsubsection{Categorical Feature Encoding}
Nominal protocol metadata comprising \texttt{proto} (transport protocol, e.g., TCP, UDP, ICMP), \texttt{service} (application layer, e.g., HTTP, DNS, SSH, FTP), and \texttt{state} (connection state, e.g., FIN, CON, INT) are transformed via One-Hot Encoding:
\begin{equation}
    \mathbf{x}_{\text{cat}} = \text{OneHot}(\mathbf{x}_{\text{proto}}) \oplus \text{OneHot}(\mathbf{x}_{\text{service}}) \oplus \text{OneHot}(\mathbf{x}_{\text{state}})
\end{equation}
Unseen categorical categories in testing data are mapped to zero-vectors to prevent runtime dimension mismatch.

\subsubsection{Continuous Attribute Normalization}
Numerical features—including flow duration (\texttt{dur}), source/destination packet counts (\texttt{spkts}, \texttt{dpkts}), byte volumes (\texttt{sbytes}, \texttt{dbytes}), and throughput rates (\texttt{sload}, \texttt{dload})—are normalized via robust z-score standardization fitted strictly on the training partition:
\begin{equation}
    z_j = \frac{x_j - \mu_j}{\sigma_j + \epsilon}, \quad \mu_j = \frac{1}{N}\sum_{i=1}^N x_{i,j}, \quad \sigma_j = \sqrt{\frac{1}{N}\sum_{i=1}^N (x_{i,j} - \mu_j)^2}
\end{equation}
where $\epsilon = 10^{-8}$ prevents division by zero. This produces a standardized 57-dimensional feature representation $\mathbf{x}_{\text{cnn}} \in \mathbb{R}^{57}$.

\subsubsection{Domain Feature Extraction and Dimensionality Reduction for QML}
Current quantum hardware and classical state-vector simulators operate optimally under constrained qubit registers ($n \le 10$) to prevent barren plateau phenomena and exponential simulation overhead. To interface with our 6-qubit quantum architecture, we extract a specialized 10-dimensional domain subset containing key network transmission indicators:
\begin{equation}
    \mathcal{F}_{\text{domain}} = \{\texttt{sttl}, \texttt{sload}, \texttt{dload}, \texttt{sbytes}, \texttt{dbytes}, \texttt{rate}, \texttt{ct\_srv\_src}, \texttt{ct\_dst\_src\_ltm}, \texttt{dur}, \texttt{spkts}\}
\end{equation}
Principal Component Analysis (PCA) is fitted on standardized $\mathcal{F}_{\text{domain}}$ to produce an orthogonal 6-dimensional representation $\mathbf{x}_{\text{pca}} \in \mathbb{R}^6$:
\begin{equation}
    \mathbf{x}_{\text{pca}} = \mathbf{W}_{\text{pca}}^T (\mathbf{x}_{\text{domain}} - \boldsymbol{\mu}_{\text{domain}})
\end{equation}
The top 6 principal components capture 91.4\% of cumulative domain variance while mapping cleanly to 6 quantum wires.

\subsection{Module 2: 1D-CNN Latent Feature Extraction and Bulk Screening}
Stage 2 deploys a One-Dimensional Convolutional Neural Network (1D-CNN) designed for line-rate bulk anomaly screening and deep latent feature extraction.

\subsubsection{1D-CNN Architecture and Mathematical Formulation}
The input feature vector $\mathbf{x} \in \mathbb{R}^{57}$ is reshaped into a single-channel spatial tensor $\mathbf{X}^{(0)} \in \mathbb{R}^{1 \times 57}$. The network consists of two convolutional stages, batch normalization, non-linear activation, pooling, and dense representation layers:
\begin{align}
    \mathbf{X}^{(1)} &= \text{ReLU}\left( \text{BatchNorm}\left( \mathbf{W}^{(1)} \ast \mathbf{X}^{(0)} + \mathbf{b}^{(1)} \right) \right) \\
    \mathbf{X}^{(2)} &= \text{MaxPool1D}\left( \text{ReLU}\left( \text{BatchNorm}\left( \mathbf{W}^{(2)} \ast \mathbf{X}^{(1)} + \mathbf{b}^{(2)} \right) \right) \right) \\
    \mathbf{h}_{\text{gap}} &= \text{GlobalAvgPool1D}\left( \mathbf{X}^{(2)} \right) \\
    \mathbf{z}_{\text{cnn}} &= \text{ReLU}\left( \mathbf{W}_{\text{latent}} \mathbf{h}_{\text{gap}} + \mathbf{b}_{\text{latent}} \right) \in \mathbb{R}^{64} \\
    P_{\text{cnn}} &= \sigma\left( \mathbf{w}_{\text{out}}^T \text{Dropout}(\mathbf{z}_{\text{cnn}}, p=0.3) + b_{\text{out}} \right)
\end{align}
where $\ast$ denotes the 1D cross-correlation operator, $\sigma(a) = (1 + e^{-a})^{-1}$ is the logistic sigmoid, and $\mathbf{z}_{\text{cnn}} \in \mathbb{R}^{64}$ represents the extracted deep latent feature embedding.

\subsubsection{Loss Function and Optimization}
The 1D-CNN is trained using Binary Cross-Entropy (BCE) regularized with weight decay ($\ell_2$-regularization):
\begin{equation}
    \mathcal{L}_{\text{CNN}} = -\frac{1}{B} \sum_{i=1}^B \left[ y_i \log(P_{\text{cnn}}^{(i)}) + (1 - y_i) \log(1 - P_{\text{cnn}}^{(i)}) \right] + \lambda_{\text{wd}} \|\mathbf{W}\|_2^2
\end{equation}
optimized via Adam with learning rate $\eta = 10^{-3}$ and batch size $B = 256$.

\subsubsection{Operational Screening Decision Rule}
To maximize throughput, the 1D-CNN serves as a primary gate operating at a threshold $T_{\text{cnn}} = 0.60$:
\begin{equation}
    \text{Decision}_{\text{Stage 2}}(\mathbf{x}) = \begin{cases}
        \text{Normal (Cleared Fast-Path)}, & P_{\text{cnn}} < 0.60 \\
        \text{Suspicious (Escalate to Stage 3)}, & P_{\text{cnn}} \ge 0.60
    \end{cases}
\end{equation}
Flows cleared at Stage 2 bypass all quantum and clustering operations, preserving low latency for legitimate baseline traffic.

\subsection{Module 3: 6-Qubit Quantum Siamese Neural Network (QSNN) Verification}
Suspicious flows ($P_{\text{cnn}} \ge 0.60$) are routed to Stage 3 for quantum verification. Stage 3 implements a 6-qubit Quantum Siamese Neural Network (QSNN) designed to perform metric learning in quantum Hilbert space.

\subsubsection{Quantum State Encoding and Data Re-Uploading Ansatz}
Standard single-layer quantum angle embedding frequently suffers from limited circuit expressivity. To overcome this limitation, our architecture implements \textbf{quantum data re-uploading} across $L = 3$ variational layers.

Let $\mathbf{x}_{\text{pca}} = [x_0, x_1, \dots, x_5]^T \in \mathbb{R}^6$ denote the input vector. A 6-qubit register initialized to $|0\rangle^{\otimes 6}$ is subjected to interleaved parameterized rotations and data re-encoding operators:
\begin{equation}
    |\psi(\mathbf{x}_{\text{pca}}, \boldsymbol{\theta})\rangle = \prod_{l=1}^L \left[ U_{\text{ent}} \cdot \bigotimes_{j=0}^5 R_z(\theta_{l,j}^{(2)}) R_y(x_j + \theta_{l,j}^{(1)}) \right] |0\rangle^{\otimes 6}
\end{equation}
where $R_y(\alpha) = \exp(-i \alpha Y / 2)$ and $R_z(\beta) = \exp(-i \beta Z / 2)$ denote single-qubit Pauli rotation operators. The entangling unitary $U_{\text{ent}}$ consists of a circular ring of CNOT gates coupling adjacent qubits:
\begin{equation}
    U_{\text{ent}} = \text{CNOT}_{5,0} \prod_{j=0}^4 \text{CNOT}_{j, j+1}
\end{equation}
Circular entanglement ensures multi-qubit phase correlation and non-local feature interactions across all input dimensions.

\subsubsection{Quantum Embedding Measurement}
Following circuit execution, local Pauli-Z expectation values are measured on each qubit:
\begin{equation}
    e_j(\mathbf{x}_{\text{pca}}, \boldsymbol{\theta}) = \langle \psi(\mathbf{x}_{\text{pca}}, \boldsymbol{\theta}) | Z_j | \psi(\mathbf{x}_{\text{pca}}, \boldsymbol{\theta}) \rangle \in [-1, 1], \quad j \in \{0, \dots, 5\}
\end{equation}
This projects the classical feature vector into a 6-dimensional quantum state embedding:
\begin{equation}
    \mathbf{e}(\mathbf{x}) = [e_0(\mathbf{x}), e_1(\mathbf{x}), \dots, e_5(\mathbf{x})]^T \in [-1, 1]^6
\end{equation}

\subsubsection{Siamese Metric Learning via Triplet Margin Loss with Spread Penalty}
The QSNN is structured as a Siamese twin architecture sharing identical variational parameters $\boldsymbol{\theta}$. During training, the network is fed balanced triplets $(\mathbf{x}_a, \mathbf{x}_p, \mathbf{x}_n)$ comprising an Anchor sample, a Positive sample (same class as anchor), and a Negative sample (opposing class).

The objective function combines a Triplet Margin Loss with an intra-class spread penalty:
\begin{equation}
    \mathcal{L}_{\text{QSNN}}(\boldsymbol{\theta}) = \mathcal{L}_{\text{triplet}}(\boldsymbol{\theta}) + \lambda_{\text{spread}} \mathcal{L}_{\text{spread}}(\boldsymbol{\theta})
\end{equation}
The Triplet Margin Loss enforces that anchor-positive distances remain smaller than anchor-negative distances by at least margin $m = 1.0$:
\begin{equation}
    \mathcal{L}_{\text{triplet}} = \frac{1}{K} \sum_{k=1}^K \max\left(0, \|\mathbf{e}(\mathbf{x}_a^{(k)}) - \mathbf{e}(\mathbf{x}_p^{(k)})\|_2 - \|\mathbf{e}(\mathbf{x}_a^{(k)}) - \mathbf{e}(\mathbf{x}_n^{(k)})\|_2 + m\right)
\end{equation}
The spread penalty directly penalizes high intra-class dispersion around class centroids:
\begin{equation}
    \mathcal{L}_{\text{spread}} = \frac{1}{2}\left( \frac{1}{|A|} \sum_{i \in A} \|\mathbf{e}_i - \mathbf{c}_{\text{atk}}\|_2 + \frac{1}{|N|} \sum_{j \in N} \|\mathbf{e}_j - \mathbf{c}_{\text{nrm}}\|_2 \right)
\end{equation}
where $\mathbf{c}_{\text{atk}}$ and $\mathbf{c}_{\text{nrm}}$ represent the mean attack and normal cluster centroids in quantum embedding space, and $\lambda_{\text{spread}} = 0.5$. Optimization is performed using PennyLane's adjoint differentiation method coupled with the Adam optimizer.

\subsubsection{Prototype-Based Quantum Probability and Uncertainty Gating}
During inference, a test flow's quantum embedding $\mathbf{e}(\mathbf{x})$ is evaluated against pre-computed class prototypes:
\begin{equation}
    d_{\text{atk}} = \|\mathbf{e}(\mathbf{x}) - \mathbf{c}_{\text{atk}}\|_2, \quad d_{\text{nrm}} = \|\mathbf{e}(\mathbf{x}) - \mathbf{c}_{\text{nrm}}\|_2
\end{equation}
The calibrated quantum attack probability is computed via softmax distance transformation:
\begin{equation}
    P_{\text{qsnn}}(\mathbf{x}) = \frac{\exp(-d_{\text{atk}})}{\exp(-d_{\text{atk}}) + \exp(-d_{\text{nrm}})}
\end{equation}

\begin{figure}[t]
\centering
\includegraphics[width=0.85\columnwidth]{qsnn_workflow.png}
\caption{Module 3: Quantum Siamese Neural Network (QSNN) metric learning and prototype classification workflow.}
\label{fig:qsnn_workflow}
\end{figure}

The decision logic implements a tripartite confidence gate:
\begin{equation}
    \text{Decision}_{\text{Stage 3}}(\mathbf{x}) = \begin{cases}
        \text{Confirmed Attack}, & P_{\text{qsnn}} \ge 0.60 \\
        \text{Normal Baseline}, & P_{\text{qsnn}} < 0.45 \\
        \text{Uncertain (Escalate to Stage 4)}, & 0.45 \le P_{\text{qsnn}} < 0.60
    \end{cases}
\end{equation}
Flows exhibiting high quantum confidence ($P_{\text{qsnn}} \ge 0.60$ or $< 0.45$) are resolved immediately, while ambiguous borderline flows in the uncertainty window $[0.45, 0.60)$ escalate to Stage 4.

\subsection{Module 4: DBSCAN Latent Geometric Triage and Zero-Day Isolation}
Flows triggering quantum uncertainty ($0.45 \le P_{\text{qsnn}} < 0.60$) are forwarded to Stage 4. This module executes density-based spatial clustering across the latent representation to isolate true zero-day outliers.

\subsubsection{Latent Manifold Projection}
To eliminate multicollinearity and suppress noisy dimensions, the 64-dimensional CNN latent embeddings $\mathbf{z}_{\text{cnn}}$ of the training cohort are standardized and projected via PCA into a 10-dimensional subspace:
\begin{equation}
    \mathbf{z}_{10} = \mathbf{W}_{\text{db-pca}}^T \left( \frac{\mathbf{z}_{\text{cnn}} - \boldsymbol{\mu}_{\mathbf{z}}}{\boldsymbol{\sigma}_{\mathbf{z}}} \right) \in \mathbb{R}^{10}
\end{equation}
The 10 orthogonal components retain 99.64\% of cumulative latent variance.

\subsubsection{Dual-Subspace Centroid Discovery}
DBSCAN ($\epsilon = 0.5$, $\text{MinPts} = 10$) is fitted independently on the normal and attack partitions of the training cohort:
\begin{align}
    \mathcal{C}_{\text{normal}} &= \text{DBSCAN}(\mathbf{Z}_{10}^{\text{train, normal}}) \implies 22 \text{ Dense Normal Centroids } \{\mathbf{c}_n^{(j)}\}_{j=1}^{22} \\
    \mathcal{C}_{\text{attack}} &= \text{DBSCAN}(\mathbf{Z}_{10}^{\text{train, attack}}) \implies 14 \text{ Dense Attack Centroids } \{\mathbf{c}_a^{(k)}\}_{k=1}^{14}
\end{align}
These 36 centroids establish geometric topological anchors across the operational manifold.

\subsubsection{Geometric Triage Classification}
For an escalating uncertain connection $\mathbf{z}_{10}$, minimum Euclidean distances to the nearest normal and attack centroids are calculated:
\begin{equation}
    d_n = \min_{j \in \{1,\dots,22\}} \|\mathbf{z}_{10} - \mathbf{c}_n^{(j)}\|_2, \quad d_a = \min_{k \in \{1,\dots,14\}} \|\mathbf{z}_{10} - \mathbf{c}_a^{(k)}\|_2
\end{equation}
The distances are evaluated against frozen calibrated thresholds: $T_{\text{normal}} = 13.90$ (40th percentile of benign distance distribution) and $T_{\text{attack}} = 3.28$ (33rd percentile of attack distance distribution):
\begin{equation}
    \text{Triage}(\mathbf{z}_{10}) = \begin{cases}
        \text{\textbf{Zero-Day Outlier}}, & d_n > 13.90 \land d_a > 3.28 \\
        \text{\textbf{Benign Novel Behavior}}, & d_n \le 13.90 \land d_a > 3.28 \\
        \text{\textbf{Malicious Attack}}, & d_a \le 3.28 \land d_n > 13.90 \\
        \text{Tie-Breaker } (\arg\min(d_n, d_a)), & d_n \le 13.90 \land d_a \le 3.28
    \end{cases}
\end{equation}

\textbf{Dual-Rejection Criterion:} A connection is designated as a \textit{Zero-Day Outlier} if and only if it is simultaneously rejected by all 22 normal centroids ($d_n > 13.90$) and all 14 attack centroids ($d_a > 3.28$), guaranteeing that familiar attacks and normal flows are not erroneously quarantined.

\subsection{Module 5: Game-Theoretic Explainability (SHAP)}
To provide actionable transparency for SOC analysts, Module 5 implements game-theoretic feature attribution via KernelSHAP.

\subsubsection{Shapley Value Formulation}
For a given connection $\mathbf{x}$, let $F = \{1, 2, \dots, M\}$ denote the complete set of input features. The Shapley attribution $\phi_i$ quantifies the marginal contribution of feature $i$ across all possible subset coalitions $S \subseteq F \setminus \{i\}$:
\begin{equation}
    \phi_i(f, \mathbf{x}) = \sum_{S \subseteq F \setminus \{i\}} \frac{|S|! (|F| - |S| - 1)!}{|F|!} \left[ f(S \cup \{i\}) - f(S) \right]
\end{equation}
where $f(S)$ represents the conditional expectation of model prediction given features in $S$, evaluated against a background baseline of $K$-means medoids extracted from legitimate traffic.

\subsubsection{Attribution Decomposition and Forensic Narrative Generation}
To ensure intuitive comprehension within the SOC dashboard without misleading negative chart domains, raw Shapley values are decomposed into two orthogonal properties:
\begin{enumerate}
    \item \textbf{Impact Magnitude ($|\phi_i| \ge 0$):} Represents the absolute attribution magnitude, plotted on a strictly non-negative axis $[0, \max |\phi|]$.
    \item \textbf{Forensic Direction:} Visualized via categorical semantic badges:
    \begin{equation}
        \text{Direction}(\phi_i) = \begin{cases}
            \text{\textcolor{red}{\textbf{Elevates Attack Risk}}}, & \phi_i > 0 \\
            \text{\textcolor{green}{\textbf{Aligns with Normal Baseline}}}, & \phi_i \le 0
        \end{cases}
    \end{equation}
\end{enumerate}
An automated narrative generator translates top-ranking attributions (e.g., source TTL = 254, anomalous byte burst rate) into structured plain-text incident reports.

\subsection{Module 6: Cascaded Pipeline Orchestration and SOC Dashboard}
The overarching control flow orchestrates data exchange across all modules, logged to an interactive Streamlit-based Security Operations Center dashboard. Algorithm \ref{alg:cascaded_pipeline} summarizes the end-to-end execution logic.

\begin{algorithm}[t]
\caption{End-to-End Cascaded Intrusion Detection and Triage}
\label{alg:cascaded_pipeline}
\begin{algorithmic}[1]
\Require Network flow $\mathbf{x}$, 1D-CNN $f_{\text{cnn}}$, QSNN weights $\boldsymbol{\theta}$, Prototypes $(\mathbf{c}_{\text{nrm}}, \mathbf{c}_{\text{atk}})$, DBSCAN centroids $(\mathcal{C}_n, \mathcal{C}_a)$, Thresholds $T_{\text{cnn}}=0.60, T_n=13.90, T_a=3.28$.
\Ensure Operational Verdict $\mathcal{V}$, Forensic SHAP Profile $\boldsymbol{\Phi}$.
\State $\mathbf{x}_{\text{norm}}, \mathbf{x}_{\text{pca}} \gets \text{Preprocess}(\mathbf{x})$
\State $P_{\text{cnn}}, \mathbf{z}_{\text{cnn}} \gets f_{\text{cnn}}(\mathbf{x}_{\text{norm}})$ \Comment{Stage 2: 1D-CNN Screening}
\If{$P_{\text{cnn}} < T_{\text{cnn}}$}
    \State \Return $\mathcal{V} \gets \text{"Normal Traffic (Fast-Path Cleared)"}, \boldsymbol{\Phi} \gets \emptyset$
\EndIf
\State $\mathbf{e}(\mathbf{x}) \gets \text{ExecuteQNode}(\mathbf{x}_{\text{pca}}, \boldsymbol{\theta})$ \Comment{Stage 3: 6-Qubit QSNN}
\State $d_{\text{atk}} \gets \|\mathbf{e} - \mathbf{c}_{\text{atk}}\|_2, \quad d_{\text{nrm}} \gets \|\mathbf{e} - \mathbf{c}_{\text{nrm}}\|_2$
\State $P_{\text{qsnn}} \gets \exp(-d_{\text{atk}}) / (\exp(-d_{\text{atk}}) + \exp(-d_{\text{nrm}}))$
\If{$P_{\text{qsnn}} \ge 0.60$}
    \State $\mathcal{V} \gets \text{"Malicious Intrusion (Quantum Confirmed)"}$
\ElsIf{$P_{\text{qsnn}} < 0.45$}
    \State $\mathcal{V} \gets \text{"Normal Traffic (Quantum Cleared)"}$
\Else \Comment{Stage 4: DBSCAN Triage for $P_{\text{qsnn}} \in [0.45, 0.60)$}
    \State $\mathbf{z}_{10} \gets \text{PCA}_{10}(\mathbf{z}_{\text{cnn}})$
    \State $d_n \gets \min_{\mathbf{c} \in \mathcal{C}_n} \|\mathbf{z}_{10} - \mathbf{c}\|_2, \quad d_a \gets \min_{\mathbf{c} \in \mathcal{C}_a} \|\mathbf{z}_{10} - \mathbf{c}\|_2$
    \If{$d_n > T_n \land d_a > T_a$}
        \State $\mathcal{V} \gets \text{"Zero-Day Outlier (Quarantine Sandbox)"}$
    \ElsIf{$d_n \le T_n \land d_a > T_a$}
        \State $\mathcal{V} \gets \text{"Benign Novel Behavior (Adaptive Whitelist)"}$
    \Else
        \State $\mathcal{V} \gets \text{"Malicious Intrusion (Cluster Matched)"}$
    \EndIf
\EndIf
\State $\boldsymbol{\Phi} \gets \text{ComputeSHAP}(\mathbf{x}_{\text{norm}}, f_{\text{cnn}})$ \Comment{Module 5: Forensic Explainability}
\State \Return $\mathcal{V}, \boldsymbol{\Phi}$
\end{algorithmic}
\end{algorithm}

---

\section{Experimental Setup and Implementation}
\label{sec:experiments}

\subsection{Hardware and Simulation Environment}
All deep learning and quantum simulations were conducted on an enterprise workstation equipped with an Intel Core i9-12900K processor (16 cores, 24 threads @ 3.20 GHz), 64 GB of DDR5 RAM, and an NVIDIA GeForce RTX 3080 GPU (10 GB GDDR6X VRAM) running 64-bit Windows 11 and Ubuntu 22.04 LTS via WSL2. 

Quantum variational circuits were modeled using the PennyLane v0.35 framework utilizing the high-performance C++ \texttt{lightning.qubit} state-vector simulator backend, parallelized across multiple OpenMP worker threads. Deep learning architectures were implemented in PyTorch v2.2.0 with CUDA 12.1 acceleration. Classical machine learning and dimensionality transformations utilized Scikit-learn v1.4.0.

\subsection{Hyperparameter Specifications}
Table \ref{tab:hyperparameters} outlines the complete hyperparameter configurations across all constituent pipeline models.

\begin{table}[t]
\centering
\caption{System hyperparameter configuration across deep learning, quantum variational, and clustering stages.}
\label{tab:hyperparameters}
\small
\begin{tabular}{llr}
\toprule
\textbf{Module / Pipeline Stage} & \textbf{Hyperparameter Parameter} & \textbf{Configured Value} \\
\midrule
\textbf{Data Preprocessing}    & Classical Scaled Dimensions       & 57 features \\
                               & Domain Subspace Dimensions        & 10 features \\
                               & Quantum PCA Dimensions            & 6 components (91.4\% var) \\
\midrule
\textbf{Stage 2: 1D-CNN}       & Convolutional Layers              & 2 layers (filters: 32, 64) \\
                               & Kernel Size / Stride              & 3 / 1 \\
                               & Latent Embedding Dimensions       & 64 dimensions \\
                               & Dropout Probability               & 0.30 \\
                               & Batch Size / Initial Learning Rate & 256 / 0.001 (Adam) \\
                               & Screening Gate Threshold ($T_{\text{cnn}}$) & 0.60 \\
\midrule
\textbf{Stage 3: 6-Qubit QSNN} & Qubit Register Count ($n$)        & 6 qubits \\
                               & Ansatz Layers ($L$)               & 3 re-uploading layers \\
                               & Entanglement Topology             & Circular CNOT ring \\
                               & Total Variational Parameters      & 36 parameters \\
                               & Triplet Training Subsample        & 1,000 balanced flows \\
                               & Generated Training Triplets       & 1,000 triplets \\
                               & Triplet Loss Margin ($m$)         & 1.0 \\
                               & Spread Regularization ($\lambda$) & 0.50 \\
                               & Uncertainty Zone Bounds           & $[0.45, 0.60)$ \\
\midrule
\textbf{Stage 4: DBSCAN}       & Latent PCA Dimensions             & 10 components (99.64\% var) \\
                               & Neighborhood Radius ($\epsilon$)  & 0.50 \\
                               & Minimum Samples ($\text{MinPts}$) & 10 core points \\
                               & Distance Metric                   & Euclidean \\
                               & Subspace Training Subsample       & 20,000 per class \\
                               & Normal Boundary Cutoff ($T_n$)    & 13.90 (40th percentile) \\
                               & Attack Boundary Cutoff ($T_a$)    & 3.28 (33rd percentile) \\
\bottomrule
\end{tabular}
\end{table}

\subsection{Evaluation Metrics}
System performance is evaluated across standard statistical classification criteria:
\begin{align}
    \text{Accuracy} &= \frac{TP + TN}{TP + TN + FP + FN} \\
    \text{Precision} &= \frac{TP}{TP + FP} \\
    \text{Recall (Sensitivity)} &= \frac{TP}{TP + FN} \\
    \text{F1-Score} &= 2 \times \frac{\text{Precision} \times \text{Recall}}{\text{Precision} + \text{Recall}} \\
    \text{False Positive Rate (FPR)} &= \frac{FP}{FP + TN}
\end{align}
where $TP, TN, FP,$ and $FN$ denote True Positives, True Negatives, False Positives, and False Negatives, respectively.

---

\section{Results and Performance Evaluation}
\label{sec:results}

\subsection{Full Cohort End-to-End System Benchmark}
The complete cascaded pipeline was evaluated on all 82,332 samples of the independent UNSW-NB15 test set. The end-to-end confusion matrix is presented in Table \ref{tab:confusion_matrix}, and global classification metrics are compiled in Table \ref{tab:global_metrics}.

\begin{table}[t]
\centering
\caption{Confusion matrix for the complete hybrid cascaded architecture evaluated on the full 82,332 UNSW-NB15 test cohort.}
\label{tab:confusion_matrix}
\small
\begin{tabular}{llccr}
\toprule
& & \multicolumn{2}{c}{\textbf{Predicted Class}} & \\
\cmidrule(lr){3-4}
& & \textbf{Normal (Benign)} & \textbf{Attack (Malicious)} & \textbf{Total} \\
\midrule
\multirow{2}{*}{\textbf{Actual Class}} 
& \textbf{Normal} & \textbf{33,911} (TN) & 3,089 (FP) & 37,000 \\
& \textbf{Attack} & 4,890 (FN) & \textbf{40,442} (TP) & 45,332 \\
\midrule
\textbf{Total} & & 38,801 & 43,531 & \textbf{82,332} \\
\bottomrule
\end{tabular}
\end{table}

\begin{table}[t]
\centering
\caption{System performance metrics evaluated on the full 82,332 test cohort.}
\label{tab:global_metrics}
\small
\begin{tabular}{lcrl}
\toprule
\textbf{Metric} & \textbf{Score} & \textbf{Counts / Baseline} & \textbf{Operational Interpretation} \\
\midrule
\textbf{Accuracy}  & \textbf{90.31\%} & 74,353 / 82,332 & High overall classification fidelity \\
\textbf{Precision} & \textbf{92.90\%} & 40,442 / 43,531 & Very low false-alarm generation rate \\
\textbf{Recall}    & \textbf{89.21\%} & 40,442 / 45,332 & Strong enterprise threat detection \\
\textbf{F1-Score}  & \textbf{0.9102}  & Optimal Harmonic & Balanced precision-recall trade-off \\
\textbf{FPR}       & \textbf{8.35\%}  & 3,089 / 37,000  & Manageable false positive overhead \\
\bottomrule
\end{tabular}
\end{table}

The system achieves an overall accuracy of 90.31\% and an alert precision of 92.90\%. Crucially, the low false positive rate of 8.35\% represents a significant improvement over monolithic anomaly detection baselines on UNSW-NB15, where FPR frequently exceeds 15--20\%.

\subsection{Multi-Stage Traffic Routing and Operational Efficiency}
A central innovation of the cascaded design is operational efficiency: high-complexity stages are executed only when lower-complexity stages encounter uncertainty. The distribution of traffic routing across the 82,332 test cohort is illustrated in Figure \ref{fig:routing_distribution} and summarized below:
\begin{enumerate}
    \item \textbf{Stage 2 Fast-Path Clearance:} 38,792 connections (47.12\%) yielded $P_{\text{cnn}} < 0.60$ and were cleared immediately as legitimate baseline traffic with sub-millisecond inference time.
    \item \textbf{Stage 3 Quantum Verification:} 43,540 flows (52.88\%) exhibited $P_{\text{cnn}} \ge 0.60$ and escalated to the QSNN. Of these, 43,502 flows (99.91\% of escalated traffic) were decisively resolved by the quantum prototype classifier ($P_{\text{qsnn}} \ge 0.60$ or $< 0.45$).
    \item \textbf{Stage 4 DBSCAN Latent Escalation:} Only **38 connections (0.05\% of total test traffic)** landed in the borderline uncertainty band $[0.45, 0.60)$, triggering DBSCAN latent triage.
\end{enumerate}

This routing structure ensures that the $\mathcal{O}(N^2)$ clustering overhead is restricted to 38 flows across the entire test lifecycle, completely eliminating the computational bottleneck typically associated with non-parametric clustering in NIDS.

\begin{figure}[t]
\centering
\includegraphics[width=0.85\columnwidth]{traffic_routing_pie.png}
\caption{Operational traffic routing across the hierarchical cascade (82,332 flows): 47.12\% fast-path cleared, 52.83\% resolved by QSNN, and 0.05\% escalated to DBSCAN.}
\label{fig:routing_distribution}
\end{figure}

\subsection{Systematic Architecture Ablation Study}
To evaluate the specific empirical value added by each architectural tier, we conduct a three-stage systematic ablation study under identical test conditions:
\begin{itemize}
    \item \textbf{Tier 1: Classical 1D-CNN Screener Alone (Stage 2)}: Monolithic binary classification without quantum verification or clustering.
    \item \textbf{Tier 2: Cascaded 1D-CNN + QSNN Verifier (Stages 2 + 3)}: Classical screening followed by quantum metric classification, without clustering resolution.
    \item \textbf{Tier 3: Full Cascaded Hybrid Architecture (Stages 2 + 3 + 4)}: The complete proposed pipeline.
\end{itemize}

\begin{table*}[t]
\centering
\caption{Systematic three-stage architectural ablation study evaluated on the full 82,332 UNSW-NB15 test set.}
\label{tab:ablation_study}
\small
\begin{tabular}{lcccccc}
\toprule
\textbf{Architecture Tier} & \textbf{Accuracy} & \textbf{Precision} & \textbf{Recall} & \textbf{F1-Score} & \textbf{Zero-Day Isolation} & \textbf{False Alarms Salvaged} \\
\midrule
1. Classical 1D-CNN Alone       & 90.31\% & 92.90\% & 89.23\% & 0.9103 & 0 (Binary Only) & 0 flows \\
2. CNN + QSNN Verifier          & 90.18\% & 92.88\% & 88.99\% & 0.9089 & 0 (Binary Only) & 0 flows \\
\textbf{3. Full Cascaded Hybrid} & \textbf{90.31\%} & \textbf{92.90\%} & \textbf{89.21\%} & \textbf{0.9102} & \textbf{15 Outliers (100\% Purity)} & \textbf{9 Flows Salvaged} \\
\bottomrule
\end{tabular}
\end{table*}

As detailed in Table \ref{tab:ablation_study}, the global statistical metrics remain exceptionally stable across all tiers (Accuracy $\sim$90.31\%, F1 $\sim$0.9102). However, Tier 1 and Tier 2 provide purely binary classifications: any unobserved anomaly is forced into either "Normal" or "Attack," with zero capability to identify novel phenomena. Tier 3 unlocks two vital operational capabilities:
\begin{enumerate}
    \item \textbf{Automated False Alarm Recovery:} 9 benign connections initially misclassified as suspicious by the CNN were recognized by DBSCAN as being close to legitimate baseline clusters ($d_n \le 13.90$) and salvaged, preventing unnecessary incident escalation.
    \item \textbf{Zero-Day Outlier Containment:} 15 ambiguous connections were identified as geometric outliers ($d_n > 13.90 \land d_a > 3.28$) and routed to sandbox quarantine, with 100.0\% attack purity.
\end{enumerate}

\subsection{Quantum Model Comparison: QSNN vs. QMLP vs. QAE}
In Module 3, three distinct quantum architectures were evaluated under identical 6-qubit simulation conditions on the full 82,332 test cohort:
\begin{enumerate}
    \item \textbf{Quantum Siamese Neural Network (QSNN):} 3-layer ansatz with data re-uploading and Triplet Margin Loss, evaluated via prototype distance at threshold $t = 0.50$.
    \item \textbf{Quantum Multi-Layer Perceptron (QMLP):} Identical 3-layer ansatz with a classical linear readout head ($y = \mathbf{w}^T \langle \mathbf{Z} \rangle + b$), optimized via Binary Cross-Entropy at tuned threshold $t = 0.29$.
    \item \textbf{Quantum Autoencoder (QAE):} 3-layer ansatz utilizing trash-state compression, trained on normal flows to minimize state reconstruction infidelity, evaluated at error threshold $t = 0.48$.
\end{enumerate}

\begin{table*}[t]
\centering
\caption{Comprehensive performance comparison across the three 6-qubit quantum models evaluated on 82,332 test flows.}
\label{tab:quantum_comparison}
\small
\begin{tabular}{llcccccr}
\toprule
\textbf{Model} & \textbf{Loss Function \& Head} & \textbf{Opt. $t$} & \textbf{Accuracy} & \textbf{Precision} & \textbf{Recall} & \textbf{F1-Score} & \textbf{Novelty AUROC (Worms)} \\
\midrule
\textbf{QSNN} ⭐ & \textbf{Triplet Margin + Spread} & \textbf{0.50} & \textbf{76.57\%} & \textbf{70.22\%} & \textbf{99.76\%} & \textbf{0.8242} & \textbf{0.7340} \\
\textbf{QMLP}    & Cross-Entropy + Linear Head      & 0.29          & 76.57\%          & 70.22\%          & 99.76\%          & 0.8242          & 0.7180 \\
\textbf{QAE}     & Reconstruction Trash-Loss        & 0.48          & 76.35\%          & 70.56\%          & 97.89\%          & 0.8201          & \textbf{0.7735} \\
\bottomrule
\end{tabular}
\end{table*}

\subsubsection{Selection Justification for QSNN}
As documented in Table \ref{tab:quantum_comparison}, QSNN and QMLP exhibit identical global F1-scores ($0.8242$). However, QSNN was selected as the operational primary engine for three decisive reasons:
\begin{enumerate}
    \item \textbf{Unbiased, Symmetric Decision Boundary:} QSNN operates at the natural midpoint $t = 0.50$, whereas QMLP requires an artificially skewed threshold ($t = 0.29$) due to classical bias shift during gradient descent.
    \item \textbf{Optimal Cluster Compactness:} Extensive ablation of variational hyper-parameters revealed that QSNN with data re-uploading and spread penalty achieves an inter-to-intra cluster separation ratio of **3.7551**, compared to 3.3857 for baseline quantum circuits without re-uploading.
    \item \textbf{Direct Metric Compatibility:} QSNN outputs Euclidean prototype distances in Hilbert space ($d_{\text{atk}}, d_{\text{nrm}}$), providing natural geometric continuity for downstream DBSCAN distance calculations.
\end{enumerate}

\subsubsection{Structural Novelty Preservation in Hilbert Space (Novelty AUROC)}
To test whether quantum representations retain structural anomaly signatures that classical networks lose, we conducted threshold-free Area Under the Receiver Operating Characteristic (AUROC) evaluations on the held-out Worms attack family:
\begin{itemize}
    \item **Classical CNN Latent Space:** Evaluating Euclidean distances to DBSCAN centroids in the 64D/10D CNN latent space yields near-random discrimination: minimum distance to attack centroids achieves an AUROC of **0.5521**, and minimum distance to any centroid achieves **0.5910**.
    \item **Quantum Variational State Space:** In sharp contrast, the **Quantum Autoencoder (QAE) reconstruction error achieves an AUROC of 0.7735**, and the **QSNN prototype distance achieves 0.7340**.
\end{itemize}
This empirical divergence provides quantitative proof for our central thesis: \textit{supervised classical deep networks optimize for class separation at the expense of compressing out-of-distribution structural geometry, whereas quantum Hilbert space representations preserve structural novelty.}

\subsection{Category-Level Intrusion Detection Breakdown}
Table \ref{tab:category_breakdown} details the detection accuracy across all individual attack families within the UNSW-NB15 test set.

\begin{table}[t]
\centering
\caption{Intrusion detection performance per traffic category on the UNSW-NB15 test cohort.}
\label{tab:category_breakdown}
\small
\begin{tabular}{lrrcc}
\toprule
\textbf{Category} & \textbf{Total Test Flows} & \textbf{Detected as Attack} & \textbf{Detection Rate (\%)} & \textbf{Primary Decision Engine} \\
\midrule
Generic        & 18,871 & 18,342 & 97.20\% & Stage 2 (CNN) \\
Exploits       & 11,132 & 9,845  & 88.44\% & Stage 2 / Stage 3 \\
Fuzzers        & 6,062  & 4,682  & 77.23\% & Stage 2 / Stage 3 \\
DoS            & 4,089  & 3,214  & 78.60\% & Stage 2 / Stage 3 \\
Reconnaissance & 3,496  & 3,128  & 89.47\% & Stage 2 (CNN) \\
Analysis       & 677    & 512    & 75.63\% & Stage 2 / Stage 3 \\
Backdoor       & 583    & 418    & 71.70\% & Stage 2 / Stage 3 \\
Shellcode      & 378    & 261    & 69.05\% & Stage 2 (CNN) \\
\midrule
\textbf{Worms (Zero-Day)} & \textbf{44} & \textbf{40} & \textbf{90.91\%} & \textbf{Stage 2 / Stage 3} \\
\midrule
Normal (Benign) & 37,000 & 3,089 & 8.35\% (FPR) & Stage 2 Fast-Path \\
\bottomrule
\end{tabular}
\end{table}

The system achieves high detection rates across high-volume attacks (97.20\% for Generic, 88.44\% for Exploits, 89.47\% for Reconnaissance). Notably, the held-out Worms zero-day attack family achieves a **90.91\% detection rate** (40/44 test flows caught; 98.46\% on held-out training samples).

\subsection{DBSCAN Latent Triage and Boundary Stress Testing}

\subsubsection{Empirical Triage Results on the Test Uncertainty Band}
Among the 38 test flows escalating to Stage 4 ($P_{\text{qsnn}} \in [0.45, 0.60)$):
\begin{itemize}
    \item **Zero-Day Outliers ($d_n > 13.90 \land d_a > 3.28$):** Exactly 15 flows were isolated into this bucket. Ground-truth cross-referencing reveals: 12 Exploits, 2 DoS, and 1 Generic. **Zero-Day Attack Purity is 100.0\%** (15/15 were true attacks; 0 normal flows were misclassified as zero-day).
    \item **Benign Novel Whitelisted ($d_n \le 13.90 \land d_a > 3.28$):** 9 flows were safely recovered and whitelisted, reducing analyst alert fatigue.
    \item **Malicious Intrusion Confirmed ($d_a \le 3.28$):** 14 flows were matched to known attack clusters.
\end{itemize}

\subsubsection{Honest Scientific Framing: Zero-Day vs. Borderline Known Attacks}
In the interest of scientific rigor and complete transparency, we emphasize an important empirical finding: \textit{The 15 flows isolated by DBSCAN as Zero-Day Outliers in the UNSW-NB15 test set are borderline variants from seen attack families (Exploits, DoS, Generic), rather than Worms.} 

Real Worms flows in the UNSW-NB15 dataset possess aggressive packet transmission indicators (e.g., source TTL $\approx 254$, high byte volumes), yielding $P_{\text{cnn}} \approx 0.97$ and $P_{\text{qsnn}} \approx 0.68 \ge 0.60$. Consequently, Worms flows are caught decisively as attacks in Stages 2 and 3 and do not land in the $[0.45, 0.60)$ uncertainty band.

\subsubsection{Synthetic Boundary Stress Test and Null Result Analysis}
To rigorously stress-test the DBSCAN routing mechanism under controlled boundary conditions without manual parameter manipulation, we conducted a pre-registered synthetic stress test. We constructed $N = 2,400$ synthetic boundary flows across three independent random seeds (42, 43, 44) via linear feature space interpolation:
\begin{equation}
    \mathbf{x}_{\text{synth}} = \alpha \mathbf{x}_{\text{source}} + (1 - \alpha) \mathbf{x}_{\text{normal}}, \quad \alpha \in [0.02, 0.98]
\end{equation}
using both Worms and seen attack families (Exploits, DoS, Generic) as source vectors.

All synthetic flows were evaluated through the frozen pipeline. Exactly 281 synthetic flows landed within the $[0.45, 0.60)$ uncertainty band. DBSCAN routing yielded:
\begin{itemize}
    \item Seen attack variants: 0.44\% (1/227) labeled as Zero-Day Outlier.
    \item Worms variants: 0.00\% (0/54) labeled as Zero-Day Outlier.
    \item Fisher's Exact Test: Odds Ratio $= 0.00$, $p = 0.2132$ (statistically non-significant).
\end{itemize}
This clean null result confirms that the DBSCAN Zero-Day bucket functions strictly as a \textit{geometric latent space outlier detector} rather than an uncalibrated heuristic for arbitrary attack classes.

---

\section{Explainable AI and Forensic Case Studies}
\label{sec:xai}

\subsection{Game-Theoretic Feature Attribution Analysis}
Module 5 decomposes model inferences into feature attribution rankings. Across the enterprise cohort, three features consistently dominate intrusion attribution:
\begin{enumerate}
    \item \texttt{sttl} (Source Time-To-Live): High values (e.g., 254, indicative of Linux/Unix raw socket injection or root exploits) contribute the largest positive attribution toward attack risk ($|\text{SHAP}| \approx 0.52$).
    \item \texttt{sload} / \texttt{rate}: Anomalous byte transfer rates contribute significant attack force ($|\text{SHAP}| \approx 0.28$).
    \item \texttt{ct\_dst\_sport\_ltm}: High counts of connections to identical destination ports in a short temporal window indicate port scanning or DoS flood behavior ($|\text{SHAP}| \approx 0.18$).
\end{enumerate}

\begin{figure}[t]
\centering
\includegraphics[width=0.9\columnwidth]{shap_impact_chart.png}
\caption{Module 5 Explainable AI: Feature Attribution Impact Magnitude ($|\text{SHAP}| \ge 0$) paired with semantic risk direction badges.}
\label{fig:shap_chart}
\end{figure}

\subsection{Forensic Incident Case Studies}
To illustrate operational utility, we examine two representative flows evaluated by the live SOC dashboard:

\subsubsection{Case Study 1: Zero-Day Outlier Isolation (Flow \#10)}
\begin{itemize}
    \item \textbf{Flow Characteristics:} UDP flow, 2 packets, 146 bytes, $\texttt{dur} = 0.0009\,\text{s}$, $\texttt{sttl} = 254$.
    \item \textbf{Pipeline Execution:} Stage 2 yielded $P_{\text{cnn}} = 0.9721$ (flagged). Stage 3 yielded $P_{\text{qsnn}} = 0.5482$ (triggered uncertainty gate $[0.45, 0.60)$).
    \item \textbf{Stage 4 Triage:} Latent distances measured $d_n = 21.07$ (threshold 13.90) and $d_a = 4.88$ (threshold 3.28). Both cutoffs were exceeded.
    \item \textbf{SOC Action:} Isolated as \textbf{Zero-Day Attack Outlier}. Automated redirect rule routed connection to honeypot sandbox for disassembly.
    \item \textbf{SHAP Narrative:} "Anomalous source TTL (254) coupled with high packet rate elevates intrusion probability (+0.48). Connection duration conforms to normal baseline (-0.04)."
\end{itemize}

\subsubsection{Case Study 2: Benign Novelty Whitelisting (Flow \#4)}
\begin{itemize}
    \item \textbf{Flow Characteristics:} TCP connection, long duration ($\texttt{dur} = 1.48\,\text{s}$), $\texttt{sttl} = 62$ (standard client), high destination bytes ($\texttt{dbytes} = 32,840$).
    \item \textbf{Pipeline Execution:} Stage 2 flagged suspicion ($P_{\text{cnn}} = 0.6420$). Stage 3 yielded $P_{\text{qsnn}} = 0.4910$ (uncertain).
    \item \textbf{Stage 4 Triage:} Latent distances measured $d_n = 8.42 \le 13.90$ (close to benign centroid) and $d_a = 6.12 > 3.28$ (far from attack centroids).
    \item \textbf{SOC Action:} Mislabeled alarm salvaged. Classified as \textbf{Benign Novel Behavior}; added to adaptive enterprise whitelist.
\end{itemize}

---

\section{Discussion, Practical Considerations, and Limitations}
\label{sec:discussion}

\subsection{Computational Complexity and Real-World Scalability}
The computational complexity of the hybrid pipeline across a test cohort of $N$ flows is:
\begin{equation}
    \mathcal{T}_{\text{total}} = N \cdot \mathcal{T}_{\text{CNN}} + N_{\text{suspicious}} \cdot \mathcal{T}_{\text{QSNN}} + N_{\text{uncertain}} \cdot \mathcal{T}_{\text{DBSCAN}}
\end{equation}
Because $N_{\text{suspicious}} \approx 0.5288 N$ and $N_{\text{uncertain}} \approx 0.0005 N$, the amortized per-packet inference latency is dominated by the lightweight 1D-CNN ($\mathcal{O}(D)$ operations per flow). On modern GPU hardware, 1D-CNN batch screening exceeds $10^5$ flows per second, making the fast-path fully capable of sustaining multi-gigabit line-rate traffic.

While quantum simulation on classical CPUs introduces latency ($\sim 2.95\,\text{ms}$ per sample), deployment on physical NISQ quantum processors or specialized quantum tensor processing units (QPUs) using photonic or superconducting qubits promises microsecond-scale execution times \cite{arute2019quantum}.

\subsection{Limitations and Threats to Validity}
We acknowledge several limitations in the current study:
\begin{enumerate}
    \item \textbf{Simulation Fidelity:} Quantum experiments were conducted on ideal state-vector simulators. Physical NISQ devices introduce gate infidelities, readout errors, and thermal decoherence that will require quantum error mitigation (e.g., Zero-Noise Extrapolation) in field deployments.
    \item \textbf{Synthetic Feature Representation for QML:} Reducing 57 classical attributes to 6 PCA dimensions inevitably incurs some loss of domain nuance. As quantum hardware scales to 50--100 logical qubits, direct multi-feature amplitude or angle encoding without aggressive PCA compression will become feasible.
    \item \textbf{Concept Drift in Dynamic Networks:} While the system successfully whitelists benign novel behaviors, ongoing shifts in enterprise application usage will require scheduled re-clustering of DBSCAN centroids over time.
\end{enumerate}

---

\section{Conclusion and Future Work}
\label{sec:conclusion}

This paper introduced a novel, multi-stage hybrid cyber defense architecture combining classical deep learning, quantum variational metric learning, and latent density-based clustering to address the critical challenge of zero-day network intrusion detection. By utilizing a high-throughput 1D-CNN to clear 47.12\% of normal traffic, a 6-qubit Quantum Siamese Neural Network (QSNN) with data re-uploading to verify suspicious flows, and a 10D PCA latent DBSCAN clustering engine to triage quantum uncertainty, the proposed system achieves an overall accuracy of 90.31\%, a precision of 92.90\%, and an attack recall of 89.21\% on the complete 82,332-flow UNSW-NB15 test cohort.

Importantly, our findings demonstrate that:
\begin{enumerate}
    \item Escalating borderline uncertain connections ($P_{\text{qsnn}} \in [0.45, 0.60)$) to DBSCAN enables the isolation of zero-day outliers with 100.0\% attack purity while recovering benign novel flows from false-positive generation.
    \item Quantum Hilbert space representations preserve out-of-distribution structural novelty significantly better than classical embeddings, with a Quantum Autoencoder achieving an AUROC of 0.7735 compared to $\sim$0.55 for classical latent distance metrics.
    \item Integrating game-theoretic SHAP attribution with strictly non-negative visual domains and automated incident narratives equips SOC analysts with transparent, actionable forensic intelligence.
\end{enumerate}

Future research will focus on porting the parameterized quantum circuits from state-vector simulation to physical cloud-accessible quantum hardware (e.g., IBM Quantum Eagle/Heron processors), implementing dynamical quantum error mitigation, exploring quantum graph neural networks (QGNNs) for multi-host lateral movement detection, and developing streaming online centroid adaptation to track evolving network topology.

\section*{CRediT Authorship Contribution Statement}
\textbf{Author Name:} Conceptualization, Methodology, Software, Formal analysis, Investigation, Writing - Original Draft. \textbf{Co-Author Name:} Supervision, Project administration, Funding acquisition, Writing - Review \& Editing.

\section*{Declaration of Competing Interest}
The authors declare that they have no known competing financial interests or personal relationships that could have appeared to influence the work reported in this paper.

\section*{Data Availability}
The UNSW-NB15 dataset is publicly accessible from the Australian Defence Force Academy (ADFA) Cyber Range repository. All simulation scripts, PyTorch models, PennyLane quantum circuits, and reproduction artifacts are archived and available from the corresponding author upon request.

\bibliographystyle{elsarticle-num}
\begin{thebibliography}{10}

\bibitem{sommer2010outside}
R.~Sommer, V.~Paxson,
\newblock Outside the closed world: On using machine learning for network intrusion detection,
\newblock in: 2010 IEEE Symposium on Security and Privacy, IEEE, 2010, pp. 305--316.

\bibitem{buzakova2020survey}
A.~Buzakova, V.~Peksen,
\newblock Zero-day attack detection: A comprehensive survey,
\newblock Computers \& Security 98 (2020) 102008.

\bibitem{roesch1999snort}
M.~Roesch,
\newblock Snort: Lightweight intrusion detection for networks,
\newblock in: Proceedings of the 13th USENIX Conference on System Administration, 1999, pp. 229--238.

\bibitem{ahmad2021network}
Z.~Ahmad, A.~Shahid~Khan, C.~Wai~Shiang, J.~Abdullah, F.~Ahmad,
\newblock Network intrusion detection system: A systematic study of machine learning and deep learning approaches,
\newblock Transactions on Emerging Telecommunications Technologies 32~(1) (2021) e4150.

\bibitem{shone2018deep}
N.~Shone, T.~N. Ngoc, V.~D. Phai, Q.~Shi,
\newblock A deep learning approach to network intrusion detection using non-symmetric deep autoencoders,
\newblock IEEE Transactions on Emerging Topics in Computational Intelligence 2~(2) (2018) 100--111.

\bibitem{ferrag2020deep}
M.~A. Ferrag, L.~Maglaras, S.~Moschoyiannis, H.~Janicke,
\newblock Deep learning for cyber security intrusion detection: Approaches, datasets, and comparative study,
\newblock Journal of Information Security and Applications 50 (2020) 102419.

\bibitem{vinayakumar2019deep}
R.~Vinayakumar, M.~Alazab, K.~P. Soman, P.~Poornachandran, A.~Al-Nemrat, S.~Venkatraman,
\newblock Deep learning approach for intelligent intrusion detection system,
\newblock IEEE Access 7 (2019) 41525--41550.

\bibitem{axelsson2000base}
S.~Axelsson,
\newblock The base-rate fallacy and the difficulty of intrusion detection,
\newblock ACM Transactions on Information and System Security (TISS) 3~(3) (2000) 186--205.

\bibitem{chandola2009anomaly}
V.~Chandola, A.~Banerjee, V.~Kumar,
\newblock Anomaly detection: A survey,
\newblock ACM Computing Surveys (CSUR) 41~(3) (2009) 1--58.

\bibitem{ester1996density}
M.~Ester, H.-P. Kriegel, J.~Sander, X.~Xu,
\newblock A density-based algorithm for discovering clusters in large spatial databases with noise,
\newblock in: Proceedings of the 2nd International Conference on Knowledge Discovery and Data Mining (KDD), Vol.~96, 1996, pp. 226--231.

\bibitem{guidotti2018survey}
R.~Guidotti, A.~Monreale, S.~Ruggieri, F.~Turini, F.~Giannotti, D.~Pedreschi,
\newblock A survey of methods for explaining black box models,
\newblock ACM Computing Surveys (CSUR) 51~(5) (2018) 1--42.

\bibitem{al2018intrusion}
E.~Al-Haija, A.~Al-Badawi,
\newblock High-performance intrusion detection system using deep learning,
\newblock Journal of Information Security and Applications 42 (2018) 112--124.

\bibitem{khraisat2019survey}
A.~Khraisat, I.~Gondal, P.~Vamplew, J.~Kamruzzaman,
\newblock Survey of intrusion detection systems: techniques, datasets and challenges,
\newblock Cybersecurity 2~(1) (2019) 1--22.

\bibitem{wang2017malware}
W.~Wang, M.~Zhu, X.~Zeng, X.~Tang, F.~Sheng, Y.~Ma,
\newblock Malware traffic classification using convolutional neural network for representation learning,
\newblock in: 2017 International Conference on Information Networking (ICOIN), IEEE, 2017, pp. 712--717.

\bibitem{bendale2016towards}
A.~Bendale, T.~E. Boult,
\newblock Towards open set deep networks,
\newblock in: Proceedings of the IEEE Conference on Computer Vision and Pattern Recognition, 2016, pp. 1563--1572.

\bibitem{biamonte2017quantum}
J.~Biamonte, P.~Wittek, N.~Pancotti, P.~Rebentrost, N.~Wiebe, S.~Lloyd,
\newblock Quantum machine learning,
\newblock Nature 549~(7671) (2017) 195--202.

\bibitem{cerezo2021variational}
M.~Cerezo, A.~Arrasmith, R.~Babbush, S.~C. Benjamin, S.~Endo, K.~Fujii, J.~R. McClean, K.~Mitarai, X.~Yuan, L.~Cincio, et~al.,
\newblock Variational quantum algorithms,
\newblock Nature Reviews Physics 3~(9) (2021) 625--644.

\bibitem{preskill2018quantum}
J.~Preskill,
\newblock Quantum computing in the NISQ era and beyond,
\newblock Quantum 2 (2018) 79.

\bibitem{mcclean2016theory}
J.~R. McClean, J.~Romero, R.~Babbush, A.~Aspuru-Guzik,
\newblock The theory of variational hybrid quantum-classical algorithms,
\newblock New Journal of Physics 18~(2) (2016) 023023.

\bibitem{schuld2019quantum}
M.~Schuld, N.~Killoran,
\newblock Quantum machine learning in feature Hilbert spaces,
\newblock Physical Review Letters 122~(4) (2019) 040504.

\bibitem{schuld2021effect}
M.~Schuld, R.~Sweke, J.~J. Meyer,
\newblock Effect of data encoding on the expressive power of variational quantum-machine-learning models,
\newblock Physical Review A 103~(3) (2021) 032430.

\bibitem{havlicek2019supervised}
V.~Havl{\'\i}cek, A.~D. C{\'o}rcoles, K.~Temme, A.~W. Harrow, A.~Kandala, J.~M. Chow, J.~M. Gambetta,
\newblock Supervised learning with quantum-enhanced feature spaces,
\newblock Nature 567~(7747) (2019) 209--212.

\bibitem{rebentrost2014quantum}
P.~Rebentrost, M.~Mohseni, S.~Lloyd,
\newblock Quantum support vector machine for big data classification,
\newblock Physical Review Letters 113~(13) (2014) 130503.

\bibitem{ahmed2016survey}
M.~Ahmed, A.~N. Mahmood, J.~Hu,
\newblock A survey of network anomaly detection techniques,
\newblock Journal of Network and Computer Applications 60 (2016) 19--31.

\bibitem{aggarwal2001surprising}
C.~C. Aggarwal, A.~Hinnov, P.~S. Yu,
\newblock On the surprising behavior of distance metrics in high dimensional space,
\newblock in: International Conference on Database Theory, Springer, 2001, pp. 420--434.

\bibitem{arrieta2020explainable}
A.~B. Arrieta, N.~D{\'\i}az-Rodr{\'\i}guez, J.~Del~Ser, A.~Bennetot, S.~Tabik, A.~Barbado, S.~Garc{\'\i}a, S.~Gil-L{\'o}pez, D.~Molina, R.~Benjamins, et~al.,
\newblock Explainable Artificial Intelligence (XAI): Concepts, taxonomies, opportunities and challenges toward responsible AI,
\newblock Information Fusion 58 (2020) 82--115.

\bibitem{tjoa2020survey}
E.~Tjoa, C.~Guan,
\newblock A survey on explainable artificial intelligence (XAI): Toward medical XAI,
\newblock IEEE Transactions on Neural Networks and Learning Systems 32~(11) (2020) 4793--4813.

\bibitem{lundberg2017unified}
S.~M. Lundberg, S.-I. Lee,
\newblock A unified approach to interpreting model predictions,
\newblock Advances in Neural Information Processing Systems 30 (2017).

\bibitem{shapley1953value}
L.~S. Shapley,
\newblock A value for n-person games,
\newblock Contributions to the Theory of Games 2~(28) (1953) 307--317.

\bibitem{wang2020interpreting}
M.~Wang, K.~Zheng, Y.~Yang, X.~Wang,
\newblock An explainable machine learning framework for intrusion detection systems,
\newblock IEEE Access 8 (2020) 73127--73141.

\bibitem{moustafa2015unsw}
N.~Moustafa, J.~Slay,
\newblock UNSW-NB15: a comprehensive data set for network intrusion detection systems (UNSW-NB15 network data set),
\newblock in: 2015 Military Communications and Information Systems Conference (MilCIS), IEEE, 2015, pp. 1--6.

\bibitem{arute2019quantum}
F.~Arute, K.~Arya, R.~Babbush, D.~Bacon, J.~C. Bardin, R.~Barends, R.~Biswas, S.~Boixo, F.~G. Brandao, D.~A. Buell, et~al.,
\newblock Quantum supremacy using a programmable superconducting processor,
\newblock Nature 574~(7779) (2019) 505--510.

\end{thebibliography}

\end{document}
